\section{What Looked Like Delivery Was Knowledge}
\label{sec:delivery}

On the hard tasks the contract beats the manual by 20 to 23 points on the
repeated panels (Section~\ref{sec:decomposition}). Both carry the same
domain, so the obvious reading, and our preprint's, is delivery: the
contract's rule is one lookup away when the query is written, while the
manual's is buried in a long prompt. The experiments in this section
replaced one difference between the arms at a time with a controlled one,
in the order the transcripts suggested it, on \mSol{}, \mLuna{} and
\mQwen{}. Their arms are named in Table~\ref{tab:arms}.

\subsection{One fact the manual leaves out}
\label{sec:one-fact}
\label{sec:interpretation}

The manual says a null fee field ``applies to all possible values''. The
fee data never uses null for its three list-typed fields; it stores ``all''
as an empty list, on 720 of 1{,}000 rules for \texttt{account\_type}. The
contract's author took that from the data and built it into the matching
predicate. \armMw{} is \armM{} plus that one fact (\emph{w}, the
wildcard), stated once in words and without SQL; its plan and the reading
of each outcome were committed before its first run. On \mSol{} it
recovers 77\% of the gap between manual and contract over all tasks
(\armM{} 65.1\%, \armMw{} 77.3\%, \armC{} 81.0\%; 90 tasks to 28 over
\armM{}), on \mQwen{} 41\%, and on \mLuna{} 38\%, as predicted for \mQwen{}
and contrary to the prediction for \mLuna{}. The fact acts where it should:
on the three families that need only the right set of fee rules, \armMw{}
reaches the contract on \mSol{} (180 of 180 task-runs against 42 for
\armM{}). On the three total-fee families it does nothing, and the
contract still wins there by 32 tasks to 0.

\subsection{A second fact, found in the transcripts}
\label{sec:second-fact}

On \mLuna{}'s 199 rule-matching tasks, those that match each payment
against the table of fee rules, \armMw{} still trailed the contract by 29
points. We read the transcripts of the 34 tasks the contract's text solved
and \armMw{} missed in both of the first two repeats. All 34 add up the
fees of many transactions, and about 30 turn on what happens when several
fee rules match one transaction, which the manual never says. The
contract's total-fee definition sums over every matched pair, the gold
convention, without marking that as an interpretation
(Section~\ref{sec:authoring-paths}). The prose arm saw the overlap and
invented a one-rule policy; twice it computed the gold number and rejected
it as double counting.

\armMws{} adds that convention to \armMw{} as a second note (\emph{s}, the
sum). In its registered run it lifts the prose on the rule-matching tasks
by 28.5 points on \mLuna{} (95 tasks to 4) and 42.7 on \mQwen{} (112 to 4).
The lift is not an artefact of where the note was found. On the 165
rule-matching tasks whose transcripts nobody read, it is 19.2 points on
\mLuna{} (64 tasks to 4) and 33.9 on \mQwen{} (78 to 4); and no \mQwen{}
transcript was read at all. \textbf{Most of what looked like delivery was
these two facts.} With both stated, the prose ties the contract's text on
\mLuna{} ($+0.8$ points, 26 tasks to 19, $p{=}0.37$).

\subsection{The facts stated alike}
\label{sec:alike}

\begin{table}[t]
\centering
\footnotesize
\setlength{\tabcolsep}{3pt}
\begin{tabular}{@{}llrr@{}}
\toprule
\tabhead{Arm} & \tabhead{Knows, and how} & \mLuna{} & \mQwen{} \\
\midrule
\armMw{} & one fact, prose & 53.9$^\dagger$ & 43.4$^\dagger$ \\
\armMws{} & both facts, prose & 84.1 & 84.1 \\
\armT{} & contract files, prompt & 80.4 & 70.7 \\
\armTws{} & files + notes, prompt & \textbf{87.8} & \textbf{89.6} \\
\armC{} & contract, governed tools & 82.6$^\dagger$ & 79.6$^\dagger$ \\
\bottomrule
\end{tabular}
\caption{Accuracy (\%) on the 199 rule-matching DABStep tasks, $k{=}3$.
The middle three arms ran together in each repeat of experiment 5c.
$^\dagger$From the earlier runs, experiment 5 (\mLuna{}) and 5b
(\mQwen{}), not rerun. The two rerun arms, \armMws{} and \armT{}, stayed
within 3 points of their earlier runs on every subset.}
\label{tab:ladder}
\end{table}

On \mQwen{}, the prose with both notes led the contract's text by 16.6
points, which no prediction foresaw. A recheck found one inequality. The
prose states each fact as a note of its own. The contract states the
summing convention less plainly: it speaks of each transaction ``joined to
\emph{the} fee rule'' before saying to sum. And its empty-list clause
comes 28{,}000 characters after a matching rule that mentions null only.
\mQwen{} often took ``the fee rule'' at its word: an approximate keyword
count finds single-rule reasoning in 103 of its 199 rule-matching runs on
the contract's text, against 37 on \mLuna{}.

\armTws{} is \armT{} plus the same two notes, word for word, and ran
together with \armMws{} and \armT{} (Table~\ref{tab:ladder}). With the
facts stated alike, the contract's text leads the prose on the
rule-matching tasks by 3.7 points on \mLuna{} (32 tasks to 13,
$p{=}0.007$) and 5.5 on \mQwen{} (35 to 10, $p{=}3{\times}10^{-4}$). That
is a small gain for the contract's form, and it is smaller than it looks.
On the 165 tasks nobody read, it is 4.6 points on \mQwen{} (23 tasks to 7,
$p{=}0.005$) but 2.2 on \mLuna{} (24 to 13, $p{=}0.10$); on \mLuna{} most
of it came from the 34 inspected tasks. And \armTws{} states each fact
twice, in the contract and in the note, so part of the lead may be
repetition. \textbf{With the facts equal, the form adds a few points at
most: clearly on \mQwen{}, not significantly on \mLuna{}'s unread tasks.}

\subsection{What the governed tools add}
\label{sec:tools-add}

\armT{} and \armC{} hold the same contract and differ only in delivery: the
files in the prompt over plain tools, or the governed lookups. On
\mLuna{}'s rule-matching tasks they tie ($-0.7$ points, 15 tasks to 18).
On its other 202 tasks the governed arm leads by 5 points (38 tasks to
15, $p{=}0.002$), most of it in three families that average fees over the
rule table (33 tasks to 4). In those families, 50 runs went right for the
governed arm and wrong for the text arm on the same task and repeat. In
every one, the text arm's answer equals, to six decimals, what one gets
by dropping the wildcard rules. The text
arm had the rule three times in its 60{,}000-character prompt. The
governed arm got it from a call made just before its first fee query: the
matching metric's lookup in 34 runs, \texttt{lookup\_domain} or the
governed \texttt{describe\_table} in 16. The validator played no part;
none of the 50 had a rejection. So delivery does buy some accuracy, by
putting the right clause in front of the agent as it writes. On \mQwen{}
the tools add 10.1 points on the rule-matching tasks (64 tasks to 16) and
trail the text on the rest (30 to 49). The effect is real, but it depends
on the model and the task family.

\subsection{Complete prose carries even a procedure}
\label{sec:procedures}

DABStep's contract gain sits on its rule-matching tasks. Two blind raters
scored the procedural load of the knowledge each question needs, on both
benchmarks shuffled together. On DABStep the gain sits on the high-load
tasks (119 tasks to 20, none on the rest), and within them on the top
level (116 to 14), where every task is the fee procedure. That left ``a
rule-matching procedure'' and ``this procedure'' confounded, so experiment
4 built new ones (Section~\ref{sec:benchmarks}): a rule-matching procedure
and a composite of dependent formulas for each of three databases, each
stated completely, once plainly and once in the style of a business manual
at twice the length. Every knowledge arm, prose or contract, in either
style, scores between 95.6\% and 100\% on each piece, on \mLuna{},
\mSol{}, and \mQwen{} (added after that ceiling was seen).
\textbf{Neither the type of knowledge nor the style of the prose decides
it}; on LiveSQLBench's own determined tasks the high-load ones gain
nothing from the contract either (8 tasks to 5, $p{=}0.58$).

\paragraph{The preprint's account, revised.} Our preprint attributed the
contract's lead over the manual to delivery: the rule one lookup away
rather than buried in a long prompt. That explanation does not hold. Most
of the lead was two facts the manual did not state. With the facts worded
alike, the contract's form adds a few points at most. The governed tools
add retrieval at the point of use, on some families and some models.
Section~\ref{sec:scale} tests burial directly: with up to four times the
base contract in the prompt, nothing is lost.
